\documentclass[12pt]{article}

\usepackage{sbc-template}
\usepackage{graphicx,url}
\usepackage[utf8]{inputenc}
\usepackage[brazil]{babel}
\usepackage{booktabs}
\usepackage{multirow}
\usepackage{amsmath}
\usepackage{array}
\usepackage{tikz}
\usetikzlibrary{shapes.geometric, arrows.meta, positioning, calc, fit, backgrounds}


\sloppy

\title{Uma  Solução Hierárquica para Detecção de Ciberataques em Ambientes IoT baseados em Redes MEC 6G }

\author{Luiz Henrique B. A. da Silva, Caio B. B. de Souza\thanks{Co-orientador}, Andson M. Balieiro\thanks{Orientador} }

\address{Curso de Sistemas de Informação\\
  Centro de Informática (CIn) -- Universidade Federal de Pernambuco (UFPE)\\
  Caixa Postal 7851 -- 50.740-560 -- Recife -- PE -- Brasil \\
  \email{lhbas@cin.ufpe.br, amb4@cin.ufpe.br, cbbs@cin.ufpe.br} 
}

\begin{document}

\maketitle

%─────────────────────────────────────────────────────────────────────────────
\begin{abstract}
The exponential growth of Internet of Things (IoT) devices and the emergence
of sixth-generation (6G) networks introduce new cybersecurity challenges in
Multi-access Edge Computing (MEC) environments. This work proposes and
implements a two-phase hierarchical Intrusion Detection System (IDS) for
IoT/MEC environments. Phase 1 runs a binary LightGBM classifier on the edge
node, optimized by the F2-score; Phase 2 classifies the specific attack type
using a multiclass LightGBM. The data
pipeline uses PCAP captures from two CIC datasets converted by
\textit{netflower}, a tool developed for this project, into a 190~GB SQLite
database. In offline evaluation, Phase 1 achieved 80\% accuracy and 96\%
attack recall (F2 = 0.886). In a fully scripted testbed validation on VIM~4
covering all eight attack categories, the binary IDS reached a 0\%
false-positive rate on a clean benign baseline (89.3\% precision and 71.1\%
recall at per-second granularity). The multiclass IDS classified
reconnaissance at F1 = 0.976; the remaining categories largely collapse onto
the DoS class within the 55 per-flow features, exposing a feature-space
limitation rather than a pipeline fault. Adding Phase~2 costs only +50~MB of
RAM with no CPU or energy overhead. Energy is reported as a calibrated
CPU-utilization estimate with an explicit uncertainty band, since the ARM node
exposes no direct power sensor. Results demonstrate the viability of the
hierarchical approach within the resource budget of an ARM edge node and
delimit where per-flow features become insufficient.
\end{abstract}

%─────────────────────────────────────────────────────────────────────────────
\begin{resumo}
O crescimento exponencial de dispositivos da Internet das Coisas (IoT) e a
emergência das redes de sexta geração (6G) introduzem novos desafios de
segurança cibernética em ambientes de Computação de Borda de Acesso Múltiplo (MEC).
Este trabalho propõe e implementa um Sistema de Detecção de Intrusão (IDS)
hierárquico de duas fases para ambientes IoT/MEC. A Fase~1 executa um
classificador binário LightGBM no nó de borda, otimizado pelo escore F2; a
Fase~2 classifica o tipo específico de ataque com LightGBM multiclasse. O pipeline utiliza PCAPs de
dois datasets do CIC, convertidos pela ferramenta \textit{netflower} em banco
SQLite de $\sim$190~GB. Em avaliação offline local, a Fase~1 atingiu acurácia de
80\% e revocação de 96\% na classe de ataque (F2 = 0,886). Em uma validação em
\textit{testbed} totalmente automatizada no VIM~4 cobrindo as oito categorias de ataque, o
IDS binário alcançou FPR de 0\% sobre um baseline benigno limpo (precisão de
89,3\% e revocação de 71,1\% por segundo). O IDS multiclasse classificou
reconhecimento com F1 = 0,976; as demais categorias colapsam majoritariamente na
classe DoS dentro das 55 \textit{features} por-fluxo, evidenciando uma limitação
do espaço de \textit{features}, não do pipeline. Adicionar a Fase~2 custa apenas
+50~MB de RAM, sem overhead de CPU ou energia. A energia é reportada como
estimativa calibrada por utilização de CPU com banda de incerteza explícita,
pois o nó ARM não expõe sensor de potência direto. Os resultados confirmam a
viabilidade da abordagem hierárquica dentro do orçamento de recursos de um nó de
borda ARM e delimitam onde as \textit{features} por-fluxo se tornam
insuficientes.
\end{resumo}

%─────────────────────────────────────────────────────────────────────────────
\section{Introdução}

A convergência entre redes móveis da Sexta Geração (6G), Internet das Coisas (IoT) e Computação de Borda de Acesso Múltiplo (MEC) promete transformar a infraestrutura
de comunicação digital com estimativas de densidades de dispositivos IoT na
ordem de $10^6$ por km² e latências de comunicação inferiores a 1~ms, sendo a tecnologia MEC um componente essencial para viabilizar esses requisitos~\cite{you:21}. Contudo,
a massificação de dispositivos expande a superfície de ataques cibernéticos, pois dispositivos IoT frequentemente possuem poder computacional e energia reduzidos, o que inviabiliza a adoção de soluções de segurança tradicionais diretamente nos terminais.

Na literatura, Sistemas de Detecção de Intrusão (IDS) baseados em regras e
assinaturas têm sido implantados em nós de borda (MEC) para aproximar a detecção
dos dispositivos IoT e reduzir a latência de resposta. Essas soluções comparam o
tráfego observado a padrões de ataque previamente catalogados e, por isso,
alcançam alta precisão sobre ameaças conhecidas~\cite{lazarevic:05}. Sua
limitação estrutural, contudo, está no que não conhecem: ataques inéditos ou
variantes de ataques existentes não correspondem a nenhuma assinatura registrada
e atravessam o detector sem gerar alerta. Além disso, manter o conjunto de
regras atualizado exige intervenção manual constante, o que se torna
impraticável na escala e na diversidade de dispositivos previstas para as redes
6G.

Diante dessas limitações, IDS baseados em aprendizado de máquina executados em
nós de borda surgem como alternativa viável: treinados com dados históricos de
tráfego, classificam fluxos de rede em tempo real e generalizam para padrões de
ataque não vistos, mantendo baixo custo computacional na
inferência~\cite{gyamfi:22}. Trabalhos recentes corroboram essa direção:
classificadores treinados sobre datasets de IoT do Canadian Institute for
Cybersecurity (CIC) alcançam alta acurácia na detecção de
ataques~\cite{houichi:25,chen:24}, arquiteturas colaborativas distribuem a
detecção entre borda e nuvem~\cite{yang:23} e esquemas de dois estágios filtram
o tráfego com um classificador leve antes de acionar um mais
custoso~\cite{aldaej:24}. Ainda assim, essas soluções enfrentam desafios em
aberto: conjuntos de dados capturados em ambientes controlados induzem desvio de
distribuição (\textit{distribution shift}) quando o modelo encontra tráfego
real, degradando a acurácia, e a diversidade de tipos de ataque demanda
objetivos de otimização distintos, como latência mínima na borda e maior
capacidade de detecção nas camadas superiores, o que motiva uma abordagem
hierárquica em que cada camada da arquitetura MEC é otimizada para um objetivo
próprio.

Este trabalho endereça essas limitações propondo um IDS
hierárquico de duas fases para ambientes IoT/MEC. A primeira fase é um
classificador binário leve, executado no nó de borda, que rotula cada fluxo de
rede como benigno ou malicioso. A segunda é um classificador multiclasse,
acionado apenas sobre os fluxos sinalizados como maliciosos, que os atribui a
classes específicas de ataque. Descreve-se o pipeline metodológico completo,
desde a coleta de dados em formato PCAP até a implantação e a avaliação da
solução em um dispositivo embarcado (Khadas VIM4).

As principais contribuições são um pipeline de dados reprodutível baseado em
SQLite para amostragem balanceada em escala, a ferramenta \textit{netflower},
publicada no PyPI, para conversão paralela de PCAPs em fluxos, o IDS hierárquico
com duas fases complementares e a validação no \textit{testbed} com tráfego de
ataque real, com comparação quantitativa entre o pipeline binário e o
multiclasse em métricas de detecção (acurácia, precisão, revocação, F1, taxas de
falsos positivos, FPR, e de falsos negativos, FNR) e de custo (consumo de CPU, RAM e energia, além de \textit{throughput} de alertas). Nos resultados mais relevantes, o IDS binário alcançou FPR de 0\% sobre um baseline benigno limpo, com precisão de 89,3\% e detecção de pelo menos 99\% dos ataques de força bruta, web, \textit{Man-in-the-Middle} (MITM) e \textit{spoofing}. A Fase~2 acrescentou apenas 50~MB de RAM, sem overhead de CPU ou energia, classificou corretamente o reconhecimento de rede (F1 = 0,976) e delimitou empiricamente onde as 55 \textit{features} por fluxo deixam de separar os tipos de ataque.

O artigo está organizado da seguinte forma. A Seção~\ref{trabalhosRelacionados} apresenta os trabalhos relacionados e a fundamentação teórica; a Seção~\ref{sec:metodologia} descreve a metodologia proposta; a Seção~\ref{sec:resultados} apresenta e discute os resultados da validação no \textit{testbed}; e a Seção~\ref{sec:conclusao} conclui o trabalho e aponta os trabalhos futuros.

%─────────────────────────────────────────────────────────────────────────────
\section{Trabalhos Relacionados}
\label{trabalhosRelacionados}

IDS podem ser classificados em baseados em assinatura e baseados em anomalia.
Os primeiros identificam padrões conhecidos com alta precisão, mas são incapazes de detectar ameaças novas. Os baseados em anomalia aprendem um perfil do comportamento normal e sinalizam desvios, sendo mais adequados para ambientes dinâmicos~\cite{lazarevic:05}. A combinação das duas abordagens em arquiteturas hierárquicas representa o estado da arte para ambientes com restrições de recursos, como dispositivos IoT e nós de borda MEC~\cite{zolanvari:19}. Esta seção revisa os trabalhos que aplicam aprendizado de máquina à detecção de intrusão em IoT, os que levam a detecção a nós MEC, os que adotam múltiplos estágios de classificação e os que tratam da segurança em redes 6G, posicionando ao final a presente proposta e as ferramentas que a fundamentam.

Os datasets do CIC são amplamente usados para treinar e avaliar IDS em redes IoT. O CIC IoT Dataset 2023~\cite{neto:23}, por exemplo, reproduz uma topologia de 105 dispositivos com 33 ataques organizados em sete categorias: Negação de Serviço Distribuída (DDoS), Negação de Serviço (DoS), \textit{recon}, ataques web, força bruta, \textit{spoofing} e Mirai. Sobre essa base, Houichi et~al.~\cite{houichi:25} comparam classificadores clássicos de aprendizado de
máquina para segurança de cidades inteligentes, enquanto Chen
et~al.~\cite{chen:24} combinam LightGBM com seleção de \textit{features} por um
algoritmo \textit{salp swarm} aprimorado para detecção de ataques em larga
escala. Entretanto, esses trabalhos apresentam apenas métricas \textit{offline} sobre o conjunto de teste do próprio dataset. Eles não avaliam o modelo implantado em
um nó de borda real, nem o desvio entre o tráfego do dataset e o de produção real, lacunas centrais endereçadas pelo trabalho aqui proposto.

Mover a detecção para a borda, por sua vez, é motivado pela baixa
latência e pelo volume de tráfego em IoT. Gyamfi e Jurcut~\cite{gyamfi:22}
revisam abordagens de projeto de IDS que combinam MEC, aprendizado de máquina e datasets, mapeando o espaço de soluções para o cenário deste
trabalho. Yang et~al.~\cite{yang:23} propõem detecção eficiente por colaboração \textit{cloud--edge}, distribuindo o processamento entre a borda e a nuvem central.
Nessas arquiteturas, porém, a etapa com maior carga de processamento é tipicamente movida para a nuvem central, e o custo de manter a detecção \textit{no próprio nó de borda}, em termos de CPU, RAM e energia, raramente é medido,  diferentemente do estudo aqui proposto. 

A ideia de filtrar o tráfego com um classificador leve antes de acionar um mais custoso também é recorrente na literatura. Aldaej et~al.~\cite{aldaej:24} propõem um IDS de dois estágios para plataformas
IoT-borda: um classificador binário inicial e, em seguida, um \textit{ensemble} para refinar a decisão sobre os fluxos sinalizados. A presente proposta compartilha essa filosofia de duas fases, mas distingue-se em dois pontos: ambas
as fases executam no mesmo nó de borda, e não distribuídas entre camadas, e o custo marginal de adicionar a segunda fase é isolado e medido diretamente no dispositivo, em vez de inferido apenas de avaliação
\textit{offline}.

No horizonte das redes 6G, esse cenário se agrava: densidade de
dispositivos e metas de latência sub-milissegundo sem precedentes~\cite{you:21}
expandem a superfície de ataque ao mesmo tempo em que reduzem a janela de
tempo disponível para detectar e reagir a uma intrusão. Sob essas restrições,
encaminhar todo o tráfego a um detector centralizado na nuvem conflita com os
requisitos de latência, tornando a inspeção no próprio nó de borda MEC um
requisito de projeto, e não mera otimização. Nessa direção, Chinnasamy
et~al.~\cite{chinnasamy:25} discutem sistemas de detecção e prevenção de intrusão
guiados por inteligência artificial para proteger redes 6G, reforçando a
necessidade de detectores que conciliem alta acurácia e baixo custo
computacional em nós de borda. É exatamente esse o regime de operação do IDS
hierárquico avaliado neste trabalho.

Esse conjunto de trabalhos avança em modelos de detecção, em arquiteturas de
borda e nos requisitos de segurança de 6G, mas nenhum deles acompanha um IDS
hierárquico do treinamento à operação em um dispositivo de borda real. É essa
lacuna que o presente trabalho ocupa, com uma demonstração ponta a ponta em que
ambas as fases executam em um nó de borda ARM de recursos limitados e a
avaliação com tráfego de ataque real complementa as métricas \textit{offline}.
O custo incremental da segunda fase, em CPU, RAM e energia, é medido diretamente
no dispositivo, e fenômenos operacionais geralmente ausentes de avaliações
puramente \textit{offline}, como o desvio de distribuição entre treino e
produção e a saturação do extrator de fluxos sob \textit{flood}, são
caracterizados e mitigados. A Tabela~\ref{tab:related} sintetiza esse
posicionamento, contrastando os trabalhos revisados segundo as cinco dimensões
que distinguem a presente proposta.

\begin{table}[ht]
\centering
\caption{Posicionamento do trabalho proposto frente à literatura segundo cinco dimensões de avaliação.}
\label{tab:related}
\setlength{\tabcolsep}{4pt}
\resizebox{\textwidth}{!}{%
\begin{tabular}{@{}lccccc@{}}
\toprule
\textbf{Trabalho}
  & \shortstack{\textbf{Deploy em}\\\textbf{borda real}}
  & \shortstack{\textbf{Avaliação}\\\textbf{online}}
  & \shortstack{\textbf{Custo no}\\\textbf{dispositivo}}
  & \shortstack{\textbf{Multi-}\\\textbf{estágio}}
  & \shortstack{\textbf{Foco}\\\textbf{6G}}
\\ \midrule
Houichi et~al.~\cite{houichi:25}       & Não & Não & Não & Não & Não \\
Chen et~al.~\cite{chen:24}             & Não & Não & Não & Não & Não \\
Gyamfi e Jurcut~\cite{gyamfi:22}       & \multicolumn{5}{c}{\textit{revisão de literatura}} \\
Yang et~al.~\cite{yang:23}             & Parcial & Não & Não & Sim & Não \\
Aldaej et~al.~\cite{aldaej:24}         & Sim & Não & Não & Sim & Não \\
Chinnasamy et~al.~\cite{chinnasamy:25} & Não & Não & Não & Não & Sim \\
\textbf{Este trabalho}                 & \textbf{Sim} & \textbf{Sim} & \textbf{Sim} & \textbf{Sim} & \textbf{Sim} \\
\bottomrule
\end{tabular}}
\end{table}

Duas ferramentas sustentam a etapa de aprendizado da proposta. O treinamento
ocorre sobre milhões de fluxos descritos por dezenas de \textit{features}, mas a
inferência precisa caber no orçamento computacional de um nó de borda ARM, o que
pede um algoritmo que combine acurácia em dados tabulares de alta
dimensionalidade com modelos compactos e rápidos na predição. O
LightGBM~\cite{ke:17} atende a esse perfil: é um método de \textit{gradient
boosting} sobre árvores de decisão com crescimento folha a folha
(\textit{leaf-wise}), que converge rapidamente e produz modelos leves o
suficiente para inferência em dispositivos restritos. Resta o ajuste de
hiperparâmetros, cuja busca exaustiva seria proibitiva no volume de dados
utilizado; o Optuna~\cite{akiba:19} trata esse problema com otimização baseada
em TPE (\textit{Tree-structured Parzen Estimator}) e poda antecipada de
tentativas pouco promissoras, reduzindo o custo de encontrar boas configurações.
Ambas as ferramentas são empregadas no treinamento das duas fases
(Seção~\ref{subsec:treinamento}).

%─────────────────────────────────────────────────────────────────────────────
\section{Metodologia}
\label{sec:metodologia}

A metodologia adotada cobre todo o ciclo do sistema, desde a coleta de capturas
brutas de rede até a avaliação do detector em operação em um nó de borda real.
A Figura~\ref{fig:metodologia} apresenta essa visão geral e o encadeamento entre
as etapas. O tráfego é coletado em arquivos de captura de pacotes (\textit{Packet
Capture}, PCAP), convertido em fluxos bidirecionais pela ferramenta
\textit{netflower} e consolidado em um banco SQLite. Sobre esse banco aplica-se a
engenharia de características que reduz o vetor de \textit{features}. Os dados
resultantes treinam os dois classificadores que compõem o IDS hierárquico e os
modelos treinados são implantados no nó de borda (placa VIM~4) para validação no
\textit{testbed} com tráfego de ataque real. As subseções seguintes detalham cada etapa na ordem da Figura~\ref{fig:metodologia}. A arquitetura hierárquica (Seção~\ref{subsec:arquitetura}), o pipeline de dados (Seção~\ref{subsec:pipeline}), o mapeamento de rótulos
(Seção~\ref{subsec:rotulos}), a engenharia de características
(Seção~\ref{subsec:features}), o treinamento das duas fases
(Seção~\ref{subsec:treinamento}) e o protocolo de validação no \textit{testbed}
(Seção~\ref{subsec:validacao}).

\begin{figure}[ht]
\centering
\resizebox{\textwidth}{!}{%
\begin{tikzpicture}[
  font=\footnotesize,
  stage/.style={rectangle, rounded corners, draw, thick, align=center,
                minimum height=1.15cm, minimum width=2.0cm, fill=gray!8},
  arr/.style={-{Stealth[length=2.5mm]}, thick}]
  \node[stage] (coleta)  {Coleta\\(PCAPs CIC +\\benigno real)};
  \node[stage, right=0.7cm of coleta] (flow)  {Extração de\\fluxos\\(\textit{netflower})};
  \node[stage, right=0.7cm of flow]   (db)    {Banco\\SQLite\\($\sim$190~GB)};
  \node[stage, right=0.7cm of db]     (feat)  {Engenharia de\\características\\(83\,$\rightarrow$\,55)};
  \node[stage, right=0.7cm of feat]   (train) {Treinamento\\(Optuna)\\Fase~1 + Fase~2};
  \node[stage, right=0.7cm of train]  (deploy){Implantação\\no nó de borda\\(VIM~4)};
  \node[stage, right=0.7cm of deploy] (val)   {Validação no\\\textit{testbed} (ataques\\+ métricas)};
  \draw[arr] (coleta) -- (flow);
  \draw[arr] (flow)   -- (db);
  \draw[arr] (db)     -- (feat);
  \draw[arr] (feat)   -- (train);
  \draw[arr] (train)  -- (deploy);
  \draw[arr] (deploy) -- (val);
\end{tikzpicture}}
\caption{Visão geral da metodologia: da coleta de PCAPs à validação no \textit{testbed} no nó de borda.}
\label{fig:metodologia}
\end{figure}

\subsection{Arquitetura hierárquica de duas fases}
\label{subsec:arquitetura}

Devido às restrições inerentes a ambientes IoT/MEC, adotou-se uma arquitetura
hierárquica de duas fases. Como a maioria do tráfego em qualquer
rede é benigna, aplicar um classificador pesado sobre todos os fluxos
desperdiçaria energia e capacidade do nó de borda. Além disso, diferentes etapas do
pipeline têm objetivos de otimização distintos: a detecção binária privilegia a
revocação, para não deixar ataques passarem, enquanto a classificação de tipo
privilegia a precisão por classe. Um pipeline progressivo concilia essas
restrições, pois a Fase~1 filtra o tráfego benigno com custo mínimo e somente os
fluxos suspeitos avançam à Fase~2 para classificação de tipo.

A arquitetura, ilustrada na Figura~\ref{fig:arquitetura}, distribui a carga
computacional hierarquicamente em duas fases complementares:

\begin{enumerate}
  \item \textbf{Fase~1 — Classificador binário (nó de borda):} decide
    \textit{benigno vs.\ ataque} com mínimo custo computacional, atuando como um
    filtro de primeira linha sobre todo o tráfego observado na interface de rede.
  \item \textbf{Fase~2 — Classificador multiclasse (nó de borda ou host):}
    identifica o tipo específico de ataque apenas para os fluxos sinalizados pela
    Fase~1. Por ser mais pesada, é a candidata natural a descarregamento para um
    host dedicado quando o nó de borda for muito restrito.
\end{enumerate}

O encadeamento das duas fases tem efeito direto sobre o custo de operação. Cada
fluxo extraído pelo \textit{netflower} é primeiro submetido à Fase~1. Caso  a
probabilidade de ataque estimada não atinja o limiar de decisão, o fluxo é
tratado como benigno e descartado sem custo adicional, de modo que o volume
dominante de tráfego legítimo nunca chega ao segundo modelo. Apenas a fração
sinalizada como suspeita aciona a Fase~2, que lhe atribui um rótulo de tipo. Como a inferência da Fase~2 ocorre em série após a Fase~1 e sobre um número reduzido de fluxos, seu custo marginal é pequeno e amortizado pelo rastreamento de fluxos, que já domina o tempo de processamento (Seção~\ref{sec:resultados}). Neste trabalho, ambas as fases foram avaliadas executando no próprio nó de
borda, um dispositivo VIM~4, enquanto um computador pessoal (PC) atuou como
gerador de ataques e como plataforma de treinamento dos modelos. Manter a Fase~2
no nó de borda, em vez de descarregá-la em um host separado, permite medir o
custo adicional do segundo modelo diretamente no dispositivo de borda
(Seção~\ref{sec:resultados}).

\begin{figure}[ht]
\centering
\begin{tikzpicture}[
  font=\footnotesize,
  box/.style={rectangle, rounded corners, draw, thick, align=center,
              minimum height=1.0cm, minimum width=3.0cm, fill=gray!8},
  dec/.style={diamond, aspect=2, draw, thick, align=center, fill=gray!12,
              inner sep=1pt},
  term/.style={rectangle, draw, thick, align=center, minimum height=0.9cm,
               minimum width=2.4cm},
  arr/.style={-{Stealth[length=2.5mm]}, thick}]
  \node[term] (traf) {Tráfego de rede\\(\texttt{eth0})};
  \node[box, below=0.8cm of traf] (f1) {\textbf{Fase~1}\\LightGBM binário};
  \node[dec, below=0.8cm of f1] (dec) {$P_1 \geq 0{,}9$?};
  \node[term, left=1.6cm of dec] (ben) {Benigno\\(descartado)};
  \node[box, below=0.9cm of dec] (f2) {\textbf{Fase~2}\\LightGBM multiclasse};
  \node[term, below=0.8cm of f2] (tipo) {Alerta tipado\\(\texttt{dos}, \texttt{recon}, \ldots)};
  \draw[arr] (traf) -- (f1);
  \draw[arr] (f1) -- (dec);
  \draw[arr] (dec) -- node[above]{não} (ben);
  \draw[arr] (dec) -- node[right]{sim} (f2);
  \draw[arr] (f2) -- (tipo);
  \begin{scope}[on background layer]
    \node[draw, dashed, rounded corners, fit=(traf)(f1)(dec)(f2)(tipo)(ben),
          inner sep=0.35cm, label={[anchor=south, font=\small, yshift=-2pt]north:nó de borda (VIM~4)}] {};
  \end{scope}
\end{tikzpicture}
\caption{Arquitetura hierárquica de duas fases executando no nó de borda (VIM~4).}
\label{fig:arquitetura}
\end{figure}

\subsection{Pipeline de dados}
\label{subsec:pipeline}

Para treinar os modelos, foi definido um pipeline de dados que converte capturas
brutas de rede em uma base consultável e amostrável em escala e que, portanto,
condiciona a qualidade de toda a detecção. O pipeline articula três elementos
que operam em sequência: os \textit{datasets} de origem, que combinam tráfego de
ataque público com tráfego benigno real coletado pelo autor; a ferramenta
\textit{netflower}, responsável por converter as capturas em fluxos
bidirecionais com \textit{features} estatísticas; e o banco SQLite, que
consolida esses fluxos para consulta e amostragem.

Para compor a base, apenas datasets com PCAPs separados por categoria foram
utilizados, nos quais os rótulos derivam do nome das pastas (\texttt{benign/},
\texttt{ddos/}, \texttt{malware/} etc.), evitando inconsistências de esquema
entre fontes~\cite{sharafaldin:18}. Foram empregados dois datasets do
CIC: o CIC IoT Dataset
2023~\cite{neto:23}\footnote{\url{https://www.unb.ca/cic/datasets/iotdataset-2023.html}} e o
CIC IIoT Dataset 2025\footnote{\url{https://www.unb.ca/cic/datasets/iiot-dataset-2025.html}},
descritos na Tabela~\ref{tab:datasets}.

\begin{table}[ht]
\centering
\caption{Datasets CIC utilizados no pipeline de dados}
\label{tab:datasets}
\begin{tabular}{@{}ll@{}}
\toprule
\textbf{Dataset} & \textbf{Foco principal} \\ \midrule
CIC IoT Dataset 2023 & Tráfego multi-ataque em dispositivos IoT \\
CIC IIoT 2025        & Intrusões em IoT industrial \\
\bottomrule
\end{tabular}
\end{table}

Além do tráfego benigno proveniente dos datasets públicos, a classe benigna foi
complementada com tráfego benigno real capturado pelo próprio autor no ambiente experimental. Ao
longo de dezenas de horas, um terminal registrou continuamente os PCAPs do uso cotidiano de um computador pessoal, incluindo tráfego de aplicações de estudo, navegação web, \textit{streaming} de vídeo, jogos, sessões SSH e atualizações de sistema, posteriormente convertidos
em fluxos pelo \textit{netflower} e rotulados como \texttt{benign}. Esse material atua como um terceiro conjunto de tráfego benigno, com a vantagem de
refletir o comportamento real de um host em uso interativo e não tráfego sintético. A motivação é metodológica. O tráfego benigno dos datasets CIC é
gerado de forma \textit{scripted}, com padrões temporais regulares, ao passo que o
tráfego interativo real apresenta características distintas, como \textit{bursts} de SSH, \textit{keepalives} HTTP e fluxos de curta duração sobre rede local rápida. Por isso, modelos treinados exclusivamente nos datasets sintéticos tendem a classificar
erroneamente o tráfego benigno real como ataque. Misturar esse tráfego real à classe benigna reduz o desvio de distribuição entre treino e implantação.

Dada a escala dos conjuntos de dados modernos de IDS, que reúnem milhões de registros e capturas brutas da ordem de terabytes, manter os dados integralmente em memória é inviável. A estratégia de ingestão e armazenamento adotada busca, portanto, escalabilidade, reprodutibilidade e execução em hardware convencional, sem a necessidade de frameworks de computação distribuída. Cada conjunto de dados é consolidado em uma única tabela SQLite\footnote{\url{https://www.sqlite.org/}}~\cite{sqlite}, preservando a estrutura original de \textit{features} e rótulos. O SQLite foi selecionado devido à sua natureza leve, portabilidade e capacidade de realizar eficientemente agregações, contagens distintas e consultas estatísticas diretamente em disco. Durante a ingestão, as \textit{features} são automaticamente mapeadas para tipos compatíveis com o SQLite: \texttt{INTEGER} para inteiros e booleanos (dado que o SQLite não possui tipo booleano nativo), \texttt{REAL} para ponto flutuante, e \texttt{TEXT} para tipos restantes ou ambíguos. Essa estratégia de mapeamento preserva a semântica original dos dados sem afetar os objetivos analíticos do estudo. A construção do banco segue duas etapas automatizadas e reprodutíveis, cujo pipeline está disponível publicamente\footnote{\url{https://github.com/luiz-linkezio/A-Hierarchical-Solution-based-on-MEC-and-IoT-for-Cyberattack-Detection-in-6G-Networks/blob/main/notebooks/database_creation.ipynb}}:

\begin{enumerate}
  \item \textbf{PCAP $\rightarrow$ CSV (\textit{Comma-Separated Values}) via netflower:} para cada pasta de categoria
    nos datasets (e.g.\ \texttt{benign/}, \texttt{ddos/}, \texttt{malware/}), todos
    os arquivos \texttt{.pcap}/\texttt{.pcapng} são convertidos em fluxos via
    \texttt{convert\_pcap\_to\_csv(n\_jobs=8)}, que executa oito processos em
    paralelo particionando os pacotes por chave bidirecional determinística. O
    rótulo de cada fluxo é derivado automaticamente do nome da pasta, dispensando
    rotulagem manual e evitando inconsistências de esquema entre as fontes. Os fluxos
    de todos os PCAPs de uma mesma pasta são então concatenados em um único arquivo
    \texttt{merged\_<label>.csv}.
  \item \textbf{CSV $\rightarrow$ SQLite:} cada \texttt{merged\_*.csv} é inserido no
    banco em \textit{chunks} de 50.000 linhas, evitando carregar o arquivo inteiro
    em memória. Ao final, cria-se um índice sobre a coluna \texttt{label}, que
    restringe as consultas por classe ao subconjunto relevante em vez de varrer toda
    a tabela. Por agregar múltiplos pacotes em um único registro de fluxo e descartar
    o detalhe pacote a pacote das capturas brutas, a representação reduz o volume em
    disco em cerca de $5\times$: o banco resultante ocupa $\sim$190~GB, ante
    $\sim$1~TB das PCAPs originais.
\end{enumerate}

A Figura~\ref{fig:pipeline} resume esse fluxo de ingestão, das pastas de
categoria à tabela única indexada por rótulo.

\begin{figure}[ht]
\centering
\resizebox{\textwidth}{!}{%
\begin{tikzpicture}[
  font=\footnotesize,
  data/.style={rectangle, draw, thick, align=center, minimum height=1.0cm,
               minimum width=1.9cm, fill=gray!8},
  proc/.style={rectangle, rounded corners, draw, thick, align=center,
               minimum height=1.0cm, minimum width=2.0cm, fill=gray!14},
  arr/.style={-{Stealth[length=2.5mm]}, thick}]
  \node[data] (pcap) {Pastas de\\categoria\\(\texttt{.pcap})};
  \node[proc, right=0.7cm of pcap] (nf) {\textit{netflower}\\\texttt{n\_jobs=8}\\(8 \textit{workers})};
  \node[data, right=0.7cm of nf] (csv) {\texttt{merged\_}\\\texttt{<label>.csv}};
  \node[proc, right=0.7cm of csv] (chunk) {Inserção em\\\textit{chunks}\\(50k linhas)};
  \node[data, right=0.7cm of chunk] (db) {SQLite\\(tabela única,\\índice em \texttt{label})};
  \draw[arr] (pcap) -- node[above]{\scriptsize fluxos} (nf);
  \draw[arr] (nf)   -- (csv);
  \draw[arr] (csv)  -- (chunk);
  \draw[arr] (chunk)-- (db);
\end{tikzpicture}}
\caption{Pipeline de ingestão: PCAPs convertidos pelo \textit{netflower} e consolidados em banco SQLite indexado por rótulo.}
\label{fig:pipeline}
\end{figure}

O \textit{netflower} é uma
ferramenta de extração de fluxos de rede desenvolvida neste trabalho, publicada no
Python Package Index (PyPI)\footnote{\url{https://pypi.org/project/netflower/}} e
disponível em código aberto sob licença MIT\footnote{\url{https://github.com/luiz-linkezio/netflower}}~\cite{netflower}.
A partir de arquivos \texttt{.pcap}/\texttt{.pcapng} ou diretamente de uma interface
de rede ao vivo, ela reconstrói fluxos bidirecionais e produz 82~\textit{features}
compatíveis com o conjunto do CICFlowMeter\footnote{\url{https://github.com/ahlashkari/CICFlowMeter}}~\cite{cicflowmeter},
organizadas em: identidade de fluxo
(IPs e portas de origem/destino, protocolo, \textit{timestamp}), estatísticas de
duração e de taxa (bytes/s e pacotes/s nos sentidos \textit{forward},
\textit{backward} e combinado), distribuições de tamanho de pacote e de tempo
entre chegadas (média, desvio-padrão, mínimo, máximo e variância), contagens das
oito \textit{flags} do protocolo da camada de transporte TCP (FIN, SYN, RST, PSH, ACK, URG, ECE, CWR), períodos
ativo/ocioso, métricas de \textit{bulk} e de \textit{subflow} e os tamanhos
iniciais da janela TCP. Apenas fluxos IPv4 sobre TCP e UDP são extraídos. Os
demais protocolos são silenciosamente ignorados.

A ferramenta foi projetada desde o início para execução em dispositivos de
borda. O \textit{parsing} dos pacotes usa o
\textit{dpkt}\footnote{\url{https://github.com/kbandla/dpkt}}~\cite{dpkt},
biblioteca 10--25$\times$ mais rápida que o
Scapy\footnote{\url{https://scapy.net}}~\cite{scapy} em processadores ARM. Para
manter a memória constante por fluxo, as estatísticas são acumuladas pelo
algoritmo \textit{online} de Welford~\cite{welford}, que atualiza média e
variância em $O(1)$ sem armazenar listas de pacotes, enquanto uma coleta
periódica de fluxos ociosos (\texttt{gc\_interval}) limita o estado residente
mesmo sob grande número de fluxos simultâneos. A captura ao vivo acessa o
libpcap\footnote{\url{https://www.tcpdump.org}}~\cite{libpcap} via
\texttt{ctypes}, o que elimina dependências de compilação no nó de borda e exige
apenas privilégio de \texttt{root} ou a \textit{capability}
\texttt{CAP\_NET\_RAW}. Na conversão de PCAPs, a função
\texttt{convert\_pcap\_to\_csv(n\_jobs)} particiona os pacotes por chave
bidirecional determinística entre processos paralelos, garantindo que pacotes de
um mesmo fluxo sejam processados pelo mesmo \textit{worker}, com escrita em CSV
bufferizada em lote para reduzir o custo de E/S e fusão posterior em um único
arquivo. Por fim, cada fluxo é emitido apenas ao completar, seja pelo término da
conexão (TCP FIN/RST), seja por \textit{timeout} de inatividade
(\texttt{idle\_timeout}, padrão 30~s) ou de duração máxima
(\texttt{flow\_timeout}, padrão 120~s), garantindo \textit{features} completas
em toda linha; esses dois \textit{timeouts} são justamente os parâmetros que
governam a latência de detecção do IDS na captura ao vivo
(Seção~\ref{subsec:validacao}).

\subsection{Mapeamento de rótulos}
\label{subsec:rotulos}

Uma vez que os ataques DDoS e DoS visam exaurir recursos e seus fluxos de rede (alta taxa de pacotes, curta duração, anomalias de flags) são estatisticamente similares, manter duas classes separadas pode adicionar ruído sem ganho discriminativo significativo no classificador. Nesse sentido, após análise das distribuições das duas classes, realizou-se a unificação delas, reduzindo o espaço de classes total de 9 para 8 (7 de ataque + 1 benigno), simplificando o problema da Fase~2. Para isso, aplicou-se $\texttt{LABEL\_MAP} = \{\texttt{ddos} \rightarrow \texttt{dos}\}$ antes do treinamento.   


Do ponto de vista do nó alvo (VIM~4), a distinção entre DoS e DDoS é de natureza
\textit{externa} aos fluxos individuais: um ataque DoS provém de uma única fonte,
enquanto o DDoS usa múltiplas fontes distribuídas, mas os fluxos individuais
recebidos pelo alvo são estatisticamente indistinguíveis. A característica (feature) que capturaria
essa diferença (diversidade de IPs de origem) exige agregação \textit{entre} fluxos
e não está presente nas 82~features por-fluxo do \textit{netflower}. Manter as
duas classes separadas no treinamento produziria um classificador que aprenderia a
distingui-las por artefatos dos datasets, não por características do tráfego real
observado pelo nó de borda. A Tabela~\ref{tab:taxonomy} apresenta a taxonomia de fluxos resultantes. 

\begin{table}[ht]
\centering
\caption{Taxonomia de ataques utilizada no treinamento do sistema IDS}
\label{tab:taxonomy}
\begin{tabular}{@{}ll@{}}
\toprule
\textbf{Categoria} & \textbf{Descrição} \\ \midrule
\texttt{benign}     & Tráfego legítimo de rede \\
\texttt{dos}        & DoS e DDoS unificados (exaustão de recursos) \\
\texttt{malware}    & Malware e comunicação de Comando e Controle (C2) \\
\texttt{bruteforce} & Ataques de força bruta (SSH, HTTP, Telnet) \\
\texttt{mitm}       & Ataques MITM (envenenamento de ARP, \textit{Address Resolution Protocol}) \\
\texttt{web}        & Ataques a aplicações web (scanning, fuzzing) \\
\texttt{spoofing}   & Falsificação de identidade de origem em pacotes IP \\
\texttt{recon}      & Varreduras e reconhecimento de rede (port scan) \\ \bottomrule
\end{tabular}
\end{table}

\subsection{Engenharia de características}
\label{subsec:features}

Partindo de 83 colunas brutas produzidas pelo \textit{netflower}, a etapa de
limpeza removeu grupos de \textit{features} que compartilham um mesmo risco:
induzir o modelo a aprender algo diferente do padrão estatístico do tráfego.

O primeiro grupo é o de identidade de fluxo (\texttt{src\_ip},
\texttt{dst\_ip}, \texttt{src\_port}, \texttt{dst\_port} e
\texttt{timestamp}). Colunas que identificam hosts e serviços específicos fazem
o modelo memorizar endereços vistos no treino em vez de aprender padrões de
tráfego, de modo que qualquer host novo no ambiente de implantação produziria
viés de previsão; todas foram removidas. Um problema análogo aparece nas colunas
de janela TCP inicial (\texttt{init\_fwd\_win\_byts} e
\texttt{init\_bwd\_win\_byts}): o \textit{netflower} escreve o valor-sentinela
$-1$ quando o tamanho da janela do \textit{handshake} não está disponível, valor
que não ocorre no tráfego real. Mantida a \textit{feature}, o modelo aprenderia
a disparar sobre o sentinela, gerando falsos positivos em fluxos legítimos nos
quais o campo é genuinamente ausente, e por isso as duas colunas foram
descartadas.

As \textit{features} de \textit{bulk rate} (\texttt{fwd\_byts\_b\_avg} e
similares) sofrem de um desvio sistemático entre treino e implantação. Elas
medem a taxa de transferência em segmentos de tráfego contínuo, o que exige que
o extrator identifique janelas de envio ininterrupto; o \textit{netflower}, em
modo ao vivo, não implementa essa heurística e emite zero para todos os fluxos,
enquanto os datasets de treino contêm valores não nulos gerados por ferramentas
que a aplicam. Manter essas colunas tornaria o modelo dependente de um sinal que
nunca existirá em produção, tornando-as não apenas inúteis, mas potencialmente
enganosas; foram, portanto, eliminadas.

Por fim, dois filtros estatísticos completaram a limpeza. Onze
\textit{features} foram removidas por correlação alta com outra ($|\rho| >
0{,}95$, Pearson): em pares altamente correlacionados, ambas carregam
essencialmente a mesma informação, e manter as duas inflaria o peso do sinal
subjacente e o custo de inferência sem ganho discriminativo; o limiar
conservador preserva \textit{features} com sobreposição apenas parcial, que
ainda podem capturar variações independentes do tráfego. Outras três foram
descartadas por variância próxima de zero (\textit{VarianceThreshold},
$\sigma^2 < 10^{-4}$): derivavam de contadores que o \textit{netflower} emite
com valor fixo na maioria dos protocolos capturados e, ao assumirem praticamente
o mesmo valor em todos os fluxos, não contribuem para separar classes, atuando
apenas como ruído constante.

O resultado foram 55 \textit{features} numéricas finais. As duas de maior
relevância pelo ganho do LightGBM foram a taxa de bytes
por segundo (\texttt{flow\_byts\_s}) e o tempo mínimo de chegada entre pacotes (\texttt{flow\_iat\_min}),
que capturam, respectivamente, a intensidade e o padrão temporal do tráfego.

\subsection{Treinamento das Duas Fases}
\label{subsec:treinamento}

As duas fases do IDS hierárquico, o classificador binário da Fase~1 e o
multiclasse da Fase~2, foram treinadas de forma independente sobre o banco
SQLite, mas com objetivos de otimização deliberadamente distintos, coerentes com
o papel de cada uma na arquitetura hierárquica. A Fase~1, como filtro de primeira linha, é
ajustada para não deixar ataques passarem, priorizando a revocação por meio do
escore F2. A Fase~2, acionada apenas sobre fluxos já sinalizados, é ajustada para
distribuir corretamente o tipo de ataque entre as classes, priorizando o macro
F1. Em ambos os casos, os hiperparâmetros foram buscados automaticamente com o
Optuna e a amostragem partiu diretamente do SQLite, garantindo balanceamento
entre classes. As próximas subseções descrevem a configuração e os resultados
\textit{offline} de cada fase.

\subsubsection{Fase 1 — Classificador Binário}

O classificador binário da Fase~1 foi treinado sobre 615.317 amostras
benignas e 615.317 de ataque (distribuídas uniformemente entre as oito classes do
SQLite), com divisão estratificada de 80/20 entre treino e teste. Tanto os
hiperparâmetros quanto o limiar de decisão foram otimizados conjuntamente pelo
Optuna (20 tentativas, 1.800~s de \textit{timeout}), tendo o escore F2 como
objetivo. A escolha do F2 é deliberada, pois ele pondera a revocação quatro vezes
mais que a precisão, refletindo que, em um IDS (e não em um Sistema de Prevenção
de Intrusão, IPS), um falso negativo (um ataque não detectado) é mais custoso que
um falso positivo (um alarme indevido).

A Tabela~\ref{tab:phase1_offline} apresenta a avaliação no conjunto de teste
(140.712 fluxos) ao limiar ótimo de 0,157. O resultado evidencia o efeito desse
objetivo: a revocação de 0,96 na classe de ataque significa que apenas 4\% dos
ataques atravessariam o filtro de primeira linha sem gerar alerta, exatamente o
comportamento esperado da fase cuja função é não deixar ataques passarem. O
custo é uma precisão de 0,68, de modo que cerca de um em cada três alertas seria
um alarme falso; em operação contínua no nó de borda, esse volume de alarmes
sobrecarregaria tanto o operador quanto a Fase~2, acionada para cada fluxo
sinalizado. O escore F2 resultante é de 0,886. Esse compromisso é aceitável em
treino, mas elevado demais para produção; por isso, o limiar de implantação foi
posteriormente fixado em 0,9, trocando parte da revocação por uma precisão
substancialmente maior e um volume de falsos positivos próximo de zero
(Seção~\ref{sec:resultados}).

\begin{table}[ht]
\centering
\caption{Avaliação offline da Fase~1 — conjunto de teste (limiar ótimo = 0,157).}
\label{tab:phase1_offline}
\begin{tabular}{@{}lrrrr@{}}
\toprule
\textbf{Classe} & \textbf{Precisão} & \textbf{Revocação} & \textbf{F1} & \textbf{Suporte} \\ \midrule
attack  & 0,68 & \textbf{0,96} & 0,80 & 57.398 \\
benign  & \textbf{0,96} & 0,69 & 0,80 & 83.314 \\ \midrule
\textbf{Acurácia} & & & \textbf{0,80} & 140.712 \\
Macro avg & 0,82 & 0,83 & 0,80 & \\
\bottomrule
\end{tabular}
\end{table}

\subsubsection{Fase 2 — Classificador Multiclasse}

O classificador multiclasse da Fase~2 foi treinado a partir de uma
amostragem independente de até 615.317 linhas por classe extraídas diretamente do
SQLite. Para compensar a escassez da classe \textit{bruteforce}
($\sim$16~mil amostras), aplicou-se \texttt{class\_weight=balanced}, e o objetivo
de otimização do Optuna foi o macro F1, que dá peso igual a todas as classes e,
portanto, penaliza explicitamente o mau desempenho nas classes minoritárias.

A Tabela~\ref{tab:phase2_offline} reporta o desempenho por classe no conjunto de
teste (317.381 fluxos). Já em avaliação \textit{offline} emerge o padrão que se
acentuará na validação no \textit{testbed}: a classe \texttt{dos} é classificada de forma
quase perfeita (F1 = 0,98), enquanto as classes de fluxo curto e baixa amostragem
\texttt{bruteforce}, \texttt{spoofing} e \texttt{mitm} apresentam precisão baixa
(0,20, 0,28 e 0,42), revelando confusão sistemática com as demais. A revocação
relativamente alta dessas classes minoritárias (0,72 a 0,77) é efeito direto do
\texttt{class\_weight=balanced}, que força o modelo a não ignorá-las, mas o macro
F1 de 0,67 confirma que a separação entre os tipos de ataque é desigual: o sinal
por-fluxo distingue bem ataques de assinatura peculiar e mal os de comportamento
semelhante a inundação. Na prática, isso significa que o rótulo de tipo emitido
pela Fase~2 é confiável para varreduras e inundações, mas ainda não para força
bruta, \textit{spoofing} e MITM, cuja triagem segue dependendo do alerta binário
da Fase~1.

\begin{table}[ht]
\centering
\caption{Avaliação offline da Fase~2 por classe — conjunto de teste.}
\label{tab:phase2_offline}
\begin{tabular}{@{}lrrrr@{}}
\toprule
\textbf{Classe} & \textbf{Precisão} & \textbf{Revocação} & \textbf{F1} & \textbf{Suporte} \\ \midrule
benign      & 0,83 & 0,69 & 0,76 & 83.313 \\
bruteforce  & 0,20 & 0,77 & 0,31 &  2.191 \\
dos         & \textbf{0,99} & \textbf{0,96} & \textbf{0,98} & 64.582 \\
malware     & 0,90 & 0,78 & 0,84 & 81.338 \\
mitm        & 0,42 & 0,72 & 0,53 & 13.823 \\
recon       & 0,83 & 0,68 & 0,75 & 47.915 \\
spoofing    & 0,28 & 0,73 & 0,41 & 12.289 \\
web         & 0,77 & 0,84 & 0,80 & 11.930 \\ \midrule
\textbf{Acurácia} & & & \textbf{0,78} & 317.381 \\
Macro F1    & & & \textbf{0,67} & \\
\bottomrule
\end{tabular}
\end{table}

\subsection{Validação no \textit{testbed}: ambiente, ferramentas e protocolo}
\label{subsec:validacao}

A validação no \textit{testbed} envolveu dois equipamentos conectados na
mesma rede local isolada, ilustrados na Figura~\ref{fig:ambiente} e especificados
na Tabela~\ref{tab:hardware}. O computador pessoal (PC) acumulou três papéis:
treinou os modelos, orquestrou todo o experimento por SSH e gerou os ataques. A
placa VIM~4 atuou como nó de borda MEC, executando os scripts IDS e monitorando a
interface de rede \texttt{eth0}, alvo de todos os ataques. O treinamento aproveitou
a GPU RTX~5070 para aceleração via CUDA (\textit{Compute Unified Device
Architecture}) do LightGBM, e os modelos treinados foram copiados para a VIM~4 para
inferência.

A VIM~4\footnote{\url{https://www.khadas.com/vim4}} foi escolhida como nó de
borda por reunir, em um computador de placa única de baixo custo, os atributos
esperados de um nó MEC de acesso. Seu \textit{System-on-Chip} (SoC) Amlogic
A311D2 combina oito núcleos ARM em arranjo heterogêneo
(Tabela~\ref{tab:hardware}): os quatro Cortex-A73 a 2,2~GHz absorvem a carga
sustentada de captura e inferência, enquanto os quatro Cortex-A53 a 1,8~GHz
atendem às tarefas leves com consumo reduzido, mantendo a placa em um envelope
de potência de 2,5 a 9~W, uma fração do consumo de um servidor convencional e
compatível com implantação em pontos de acesso com alimentação restrita. Os
8~GB de RAM LPDDR4x dão margem para os modelos das duas fases residirem em
memória simultaneamente. Em conectividade, além da Gigabit Ethernet usada nos
experimentos, a placa oferece Wi-Fi~6, Bluetooth~5.1, USB~3.0 e GPIO de 40
pinos, o que permite agregar sensores e enlaces heterogêneos típicos de
ambientes IoT. Por fim, o suporte oficial a Ubuntu Linux e a Android amplia o
reuso do mesmo hardware em outros papéis, de \textit{gateway} IoT a plataforma
de aplicações Android embarcadas, reforçando o custo-benefício da escolha.

\begin{figure}[ht]
\centering
\definecolor{cardfill}{HTML}{F5F7FA}
\definecolor{cardstroke}{HTML}{334155}
\definecolor{devfill}{HTML}{E2E8F0}
\definecolor{devstroke}{HTML}{334155}
\definecolor{screenfill}{HTML}{93C5FD}
\definecolor{boardfill}{HTML}{CBD5E1}
\definecolor{chipfill}{HTML}{475569}
\definecolor{netblue}{HTML}{2563EB}
\definecolor{netteal}{HTML}{0F766E}
\resizebox{\textwidth}{!}{%
\begin{tikzpicture}[
  font=\footnotesize, line join=round,
  card/.style={rectangle, rounded corners=5pt, draw=cardstroke, thick,
               fill=cardfill, inner sep=9pt},
  subbox/.style={rectangle, rounded corners=3pt, draw=devstroke, semithick,
                 fill=white, align=center, font=\scriptsize, inner sep=5pt,
                 text width=2.7cm},
  logbox/.style={subbox, fill=blue!6, draw=netblue!70},
  callarr/.style={-{Stealth[length=2mm]}, thick, draw=devstroke},
  dataarr/.style={-{Stealth[length=2mm]}, thick, draw=netteal},
  attackarr/.style={-{Stealth[length=3mm]}, line width=1pt, draw=netblue,
                    dash pattern=on 0.5pt off 2.5pt},
  collectarr/.style={-{Stealth[length=2.5mm]}, line width=0.9pt, draw=netteal,
                     dash pattern=on 0.5pt off 2.5pt}]

  %======== PC: hardware + elementos de software ========
  \node (pcimg) {%
    \begin{tikzpicture}[line join=round]
      % gabinete
      \draw[thick, draw=devstroke, fill=devfill, rounded corners=1pt]
            (-0.80,-0.35) rectangle (-0.50,0.62);
      \draw[devstroke] (-0.75,0.50)--(-0.55,0.50);
      \draw[devstroke] (-0.75,0.42)--(-0.55,0.42);
      \fill[netblue] (-0.65,-0.16) circle(0.035);
      % monitor
      \draw[thick, draw=devstroke, fill=devfill, rounded corners=2pt]
            (-0.35,0.00) rectangle (0.80,0.62);
      \fill[screenfill] (-0.28,0.08) rectangle (0.73,0.54);
      \draw[thick, devstroke] (0.22,0.00)--(0.22,-0.18);
      \draw[line width=1.4pt, devstroke] (0.02,-0.20)--(0.42,-0.20);
    \end{tikzpicture}};
  \node[below=3pt of pcimg, align=center] (pctitle)
    {\textbf{PC}\\Ryzen~5 7600 + RTX~5070\\{\scriptsize treino dos modelos (GPU)}};
  \node[subbox, below=8pt of pctitle] (orch)
    {\textbf{Orquestrador}\\\texttt{\scriptsize run\_experiment.sh}\\[3pt]
     \tiny deploy SSH, sessões A/B, teardown, métricas};
  \node[subbox, below=9pt of orch] (atk)
    {\textbf{Gerador de ataques}\\\texttt{\scriptsize attack\_generator.py}\\[3pt]
     \tiny 8 categorias de ataque};
  \node[logbox, below=9pt of atk] (relat)
    {\textbf{Relatório de ataques}\\\texttt{\scriptsize report\_*.json}\\[3pt]
     \tiny janelas de \textit{ground truth} por ataque};
  \node[subbox, below=9pt of relat] (metrics)
    {\textbf{Métricas}\\\texttt{\scriptsize ids\_metrics.py}\\[3pt]
     \tiny cruza log do IDS + relatório; acurácia, F1, energia};
  \draw[callarr] (orch) -- node[right, font=\tiny, text=black, xshift=1pt]{invoca} (atk);
  \draw[dataarr] (atk) -- node[right, font=\tiny, text=black, xshift=1pt]{gera} (relat);
  \draw[dataarr] (relat) -- node[right, font=\tiny, text=black, xshift=1pt]{coleta} (metrics);

  %======== VIM 4: hardware + elementos de software ========
  \node[right=7.5cm of pcimg] (vimimg) {%
    \begin{tikzpicture}[line join=round]
      % placa (PCB)
      \draw[thick, draw=devstroke, fill=boardfill, rounded corners=2pt]
            (-0.82,-0.52) rectangle (0.82,0.52);
      \foreach \x/\y in {-0.70/0.40, 0.70/0.40, -0.70/-0.40, 0.70/-0.40}
        \draw[devstroke, fill=white] (\x,\y) circle(0.05);
      % SoC + dissipador
      \draw[thick, draw=devstroke, fill=chipfill] (-0.28,-0.28) rectangle (0.22,0.28);
      \foreach \x in {-0.22,-0.11,0.00,0.11,0.18}
        \draw[white, line width=0.6pt] (\x,-0.24)--(\x,0.24);
      % barra de pinos (GPIO)
      \foreach \x in {-0.58,-0.49,-0.40,-0.31}
        \fill[chipfill] (\x,0.36) rectangle ++(0.05,0.10);
      % portas (USB / Ethernet)
      \draw[thick, draw=devstroke, fill=devfill] (0.52,0.12) rectangle (0.86,0.34);
      \draw[thick, draw=devstroke, fill=devfill] (0.52,-0.16) rectangle (0.86,0.06);
    \end{tikzpicture}};
  \node[below=3pt of vimimg, align=center] (vimtitle)
    {\textbf{VIM~4}\\Amlogic A311D2 (ARM)\\{\scriptsize nó de borda MEC}};
  \node[subbox, below=8pt of vimtitle] (idsbox)
    {\textbf{IDS}\\\texttt{\scriptsize network\_ids.py}\\[3pt]
     \tiny Fase~1 (+ Fase~2)\\interface \texttt{eth0}};
  \node[logbox, below=9pt of idsbox] (idslog)
    {\textbf{Log do IDS}\\\texttt{\scriptsize ids\_run\_*.log}\\[3pt]
     \tiny alertas por fluxo + \textit{snapshots} de sistema};
  \draw[dataarr] (idsbox) -- node[right, font=\tiny, text=black, xshift=1pt]{registra} (idslog);

  %======== cartões de fundo ========
  \begin{scope}[on background layer]
    \node[card, fit=(pcimg)(pctitle)(orch)(atk)(relat)(metrics)] (pccard) {};
    \node[card, fit=(vimimg)(vimtitle)(idsbox)(idslog)] (vimcard) {};
  \end{scope}

  %======== switch / rede isolada (centro) ========
  \coordinate (midx) at ($(pccard.east)!0.5!(vimcard.west)$);
  \node (sw) at (midx |- idsbox) {%
    \begin{tikzpicture}[line join=round]
      \draw[thick, draw=devstroke, fill=devfill, rounded corners=1.5pt]
            (-0.52,-0.16) rectangle (0.52,0.16);
      \foreach \x in {-0.40,-0.24,-0.08,0.08,0.24}
        \draw[devstroke, fill=white] (\x,-0.10) rectangle ++(0.10,0.12);
      \fill[netblue] (0.45,0.09) circle(0.022);
    \end{tikzpicture}};
  \node[above=1pt of sw, font=\scriptsize\itshape, text=black]
    {Rede local isolada (Gigabit Ethernet)};

  %======== conexões entre os nós (rede isolada) ========
  % tráfego de ataque: sai do gerador de ataques direto para o IDS da VIM
  \draw[attackarr] (atk.east) --
        node[pos=0.5, below=3pt, align=center, font=\scriptsize, text=black]
        {ataques (8 categorias)\\\texttt{192.168.100.5}}
        (idsbox.west);
  % coleta de log: ids\_metrics.py busca o log do IDS na VIM 4 via SSH
  \draw[collectarr] (idslog.west) -- ++(-0.7,0) |- (metrics.east);
  \node[above=2pt, font=\scriptsize, text=black] at (midx |- metrics)
    {coleta do log (via SSH)};
\end{tikzpicture}}
\caption{Ambiente de experimentação: PC e VIM~4 conectados por rede local isolada.}
\label{fig:ambiente}
\end{figure}

\begin{table}[ht]
\centering
\caption{Especificações do hardware utilizado}
\label{tab:hardware}
\begin{tabular}{@{}lll@{}}
\toprule
\textbf{Componente} & \textbf{PC (treino / ataques)} & \textbf{VIM 4 (IDS)} \\ \midrule
SoC / CPU & AMD Ryzen 5 7600 (6c/12t) & Amlogic A311D2 \\
          &                            & 4$\times$ Cortex-A73 @ 2,2~GHz \\
          &                            & + 4$\times$ Cortex-A53 @ 1,8~GHz \\
GPU       & NVIDIA RTX 5070 (16~GB)    & Mali-G52 MP8 \\
RAM       & 2$\times$16~GB DDR5 5600~MHz & 8~GB LPDDR4x \\
Placa-mãe & Asus TUF Gaming B650 Plus WiFi & — \\
Fonte     & MSI A750GF (750~W, 80+ Gold) & — \\
Rede      & Gigabit Ethernet           & Gigabit Ethernet \\
\bottomrule
\end{tabular}
\end{table}

A Figura~\ref{fig:ambiente} detalha os elementos de software que executam sobre
esse hardware. No PC residem quatro elementos. O orquestrador
(\texttt{run\_experiment.sh}) coordena todo o experimento: implanta os scripts na
VIM~4 via SSH, conduz as duas sessões de avaliação (A e B, descritas adiante),
encerra os processos ao final de cada sessão e dispara o cálculo de métricas. O
gerador de ataques (\texttt{attack\_generator.py}), invocado pelo orquestrador,
executa as oito categorias de ataque em sequência contra a VIM~4. Ao término de
cada execução, o gerador produz o relatório de ataques (\texttt{report\_*.json}),
que registra as janelas temporais de início e fim de cada ataque e constitui o
\textit{ground truth} da avaliação. Por fim, o script de métricas
(\texttt{ids\_metrics.py}) coleta via SSH o log produzido pelo IDS na VIM~4 e o
cruza com o relatório de ataques para calcular as métricas de detecção, de
recursos e de energia. Na VIM~4, por sua vez, executa o IDS propriamente dito
(\texttt{network\_ids.py}), que monitora a interface \texttt{eth0}, roda a Fase~1
(e, quando habilitada, a Fase~2) e registra sua operação no log do IDS
(\texttt{ids\_run\_*.log}), contendo os alertas emitidos por fluxo e
\textit{snapshots} periódicos do estado do sistema (CPU, RAM, temperatura e
frequência), dos quais derivam as métricas de custo.

Para executar as 8 categorias de ataque, o gerador emprega as seguintes
ferramentas:

\begin{itemize}
  \item \textbf{nmap}\footnote{\url{https://nmap.org}}~\cite{nmap} — varredura de
    rede: SYN scan (\texttt{-sS}), UDP scan dos 100 top ports (\texttt{-sU}),
    varredura agressiva (\texttt{-A});
  \item \textbf{masscan}\footnote{\url{https://github.com/robertdavidgraham/masscan}}~\cite{masscan}
    — varredura SYN de alta velocidade;
  \item \textbf{hping3}\footnote{\url{http://www.hping.org}}~\cite{hping3} —
    geração de inundação: SYN flood (porta 80), UDP flood (porta 53), ICMP flood;
    DDoS com fonte aleatória (\texttt{--rand-source}) nas portas 443 e 80;
    spoofing com fonte forjada fixa (\texttt{-a});
  \item \textbf{medusa}\footnote{\url{https://github.com/jmk-foofus/medusa}}~\cite{medusa}
    — força bruta de credenciais: SSH (root), HTTP (admin), Telnet (root);
  \item \textbf{nikto}\footnote{\url{https://github.com/sullo/nikto}}~\cite{nikto}
    — varredura de vulnerabilidades web;
  \item \textbf{gobuster}\footnote{\url{https://github.com/OJ/gobuster}}~\cite{gobuster}
    — enumeração de diretórios/caminhos HTTP;
  \item \textbf{curl}\footnote{\url{https://curl.se}}~\cite{curl} — inundação HTTP
    (100 requisições em paralelo);
  \item \textbf{arpspoof} (suíte \textit{dsniff})\footnote{\url{https://www.monkey.org/~dugsong/dsniff/}}~\cite{dsniff}
    — envenenamento ARP bidirecional (ataque MITM), com tráfego forçado pela rota
    envenenada;
  \item \textbf{nc} (netcat)\footnote{\url{https://nmap.org/ncat/}}~\cite{ncat} —
    beacon C2 simulado via TCP porta~4444 (100 conexões curtas).
\end{itemize}

Os dois pipelines são avaliados em duas sessões
sequenciais sob o mesmo roteiro: a Sessão~A apenas com o IDS binário e a Sessão~B
com o binário+multiclasse, o que dá atribuição limpa de recursos a cada pipeline
e janelas de \textit{ground truth} estruturalmente idênticas. Cada sessão inclui
60~s de baseline benigno antes e depois dos ataques e 120~s de intervalo
ocioso entre ataques consecutivos, para que os fluxos de cada ataque drenem
antes da janela seguinte (motivado pela saturação do extrator sob \textit{flood},
discutida na Seção~\ref{subsec:res_artefatos}). Os ataques visam \texttt{192.168.100.5}
(VIM~4), com sincronização por \textit{offset} de 3~h de BRT para UTC; limiar
de detecção binária $P_1 \geq 0{,}9$; \textit{netflower} com \texttt{idle\_timeout}
= 30~s e \texttt{flow\_timeout} = 120~s. A captura PCAP no atacante foi
deliberadamente desativada: as métricas derivam do log do IDS (alertas + amostras
periódicas de sistema) e das janelas temporais do relatório do gerador de ataques, não
dos PCAPs, que, sob os \textit{floods}, chegavam a vários GB sem valor analítico.

A avaliação binária é feita por
segundo: cada segundo é rotulado ``ataque'' se cai em alguma janela de ataque
(estendida por \texttt{idle\_slack} = 90~s, que absorve o rastro de emissão
atrasada dos \textit{floods}) e ``benigno'' caso contrário, com o \textit{ground
truth} derivado das janelas do gerador de ataques; o FPR é medido sobre o baseline
benigno \textit{pré-ataque} (limpo). A avaliação multiclasse mapeia cada alerta à
sua janela e compara o tipo previsto ao verdadeiro. CPU, RAM e \textit{throughput}
derivam das amostras periódicas do log; a energia segue um modelo linear de CPU
calibrado, detalhado na Seção~\ref{subsec:res_comp}.

%─────────────────────────────────────────────────────────────────────────────
\section{Resultados}
\label{sec:resultados}

Esta seção apresenta os resultados da validação no \textit{testbed} no VIM~4, cujo ambiente,
organização do experimento e protocolo foram descritos na
Seção~\ref{subsec:validacao}. A exposição segue do particular ao geral: primeiro
avaliam-se as duas fases isoladamente, o IDS binário da Fase~1
(Seção~\ref{subsec:res_bin}) e o IDS multiclasse das Fases~1+2
(Seção~\ref{subsec:res_multi}). Em seguida, caracterizam-se os artefatos de
implantação que condicionam esses números (Seção~\ref{subsec:res_artefatos}). Por fim, confrontam-se os dois pipelines em uma comparação direta de desempenho de detecção e consumo de recursos, \textit{throughput} e energia (Seção~\ref{subsec:res_comp}).

Antes dos números, convém precisar a semântica das métricas reportadas e por que
são adequadas a um IDS. Nas equações a seguir, TP, FP, TN e FN denotam,
respectivamente, os verdadeiros positivos, falsos positivos, verdadeiros
negativos e falsos negativos. A precisão é a fração de alertas que correspondem
a ataques reais,
\begin{equation}
\text{Precisão} = \frac{\mathrm{TP}}{\mathrm{TP} + \mathrm{FP}},
\label{eq:precisao}
\end{equation}
de modo que uma precisão baixa significa muitos alarmes falsos, que
sobrecarregam o operador. A revocação (ou taxa de verdadeiros positivos, TPR) é
a fração de ataques efetivamente detectados,
\begin{equation}
\text{Revocação} = \frac{\mathrm{TP}}{\mathrm{TP} + \mathrm{FN}},
\label{eq:revocacao}
\end{equation}
e uma revocação baixa significa ataques que passam despercebidos. As duas se
complementam, e o F1 é sua média harmônica, um resumo único que só é alto quando
ambas o são:
\begin{equation}
\mathrm{F1} = 2 \times \frac{\text{Precisão} \times \text{Revocação}}{\text{Precisão} + \text{Revocação}}.
\label{eq:f1}
\end{equation}
A taxa de falsos positivos (\textit{False Positive Rate}, FPR) é a fração de
tráfego benigno classificado como ataque,
\begin{equation}
\mathrm{FPR} = \frac{\mathrm{FP}}{\mathrm{FP} + \mathrm{TN}},
\label{eq:fpr}
\end{equation}
que em um IDS é crítica, pois um FPR alto inviabiliza a operação ao afogar a
equipe em alarmes sobre tráfego legítimo. A taxa de falsos negativos
(\textit{False Negative Rate}, FNR) mede o complemento, os ataques perdidos:
\begin{equation}
\mathrm{FNR} = \frac{\mathrm{FN}}{\mathrm{FN} + \mathrm{TP}} = 1 - \text{Revocação}.
\label{eq:fnr}
\end{equation}
Como num IDS um ataque não detectado é mais custoso que um alarme indevido, este
trabalho privilegia revocação e baixo FPR, e usa F1 e macro F1 (média do
F1 por classe, que dá peso igual a classes raras) para avaliar o equilíbrio entre
os tipos de ataque na Fase~2.

\subsection{IDS Binário (Fase 1)}
\label{subsec:res_bin}

A Tabela~\ref{tab:bin_cm} apresenta a matriz de confusão por segundo do classificador binário (Fase 1). O número que mais importa para operar o IDS é o FPR medido sobre o baseline
benigno pré-ataque, sem tráfego de ataque anterior drenando: zero falsos
positivos em 57~s (Tabela~\ref{tab:bin_metrics}), ou seja, nenhum alarme indevido
sobre tráfego legítimo limpo. Os 116 ``FP'' da matriz por segundo
(Tabela~\ref{tab:bin_cm}) não contradizem esse valor, pois não são erros sobre
tráfego benigno: são fluxos de \textit{flood} emitidos com atraso pelo extrator
após o fim do ataque, quando o nó ARM satura sob a alta taxa de chegada
(Seção~\ref{subsec:res_artefatos}). Por isso o FPR ``global'' por segundo
(44,6\%) não é representativo da operação real.

\begin{table}[ht]
\centering
\caption{Matriz de confusão do IDS binário — granularidade por segundo}
\label{tab:bin_cm}
\begin{tabular}{@{}llrr@{}}
\toprule
& & \multicolumn{2}{c}{\textbf{Predição}} \\
\cmidrule{3-4}
& & \textbf{Ataque} & \textbf{Benigno} \\ \midrule
\multirow{2}{*}{\textbf{Ground truth}}
  & Ataque & 967 & 393 \\
  & Benigno & 116 & 144 \\ \bottomrule
\end{tabular}
\end{table}

\begin{table}[ht]
\centering
\caption{Métricas do IDS binário na validação no \textit{testbed} (limiar = 0,9).}
\label{tab:bin_metrics}
\begin{tabular}{@{}lr@{}}
\toprule
\textbf{Métrica} & \textbf{Valor} \\ \midrule
Acurácia          & 68,6\% \\
Precisão          & \textbf{89,3\%} \\
Revocação (TPR)   & 71,1\% \\
F1-score          & 79,2\% \\
\textbf{FPR (baseline benigno limpo)} & \textbf{0,0\%} (0/57~s) \\
FNR               & 28,9\% \\ \bottomrule
\end{tabular}
\end{table}



Os 393 falsos negativos por segundo (Tabela~\ref{tab:detection_rate}) não se
distribuem por igual entre os ataques. As classes de inundação e curta duração
(força bruta, web, MITM e \textit{spoofing}) atingem pelo menos 99\% de detecção
por segundo, ao passo que os ataques perdidos concentram-se em \textit{recon}
(20,9\% de detecção), \texttt{dos} (59,4\%) e \textit{malware} (21,3\%). A baixa
detecção de \textit{recon} é consequência direta do \textit{scan} UDP do nmap:
poucos pacotes por porta não acumulam \textit{features} suficientes para que o
fluxo seja emitido dentro do \texttt{idle\_timeout} de 30~s, e o ataque, embora
em curso, permanece invisível ao detector até que a janela expire. Essa é
justamente a classe de ataque que uma segunda fase de estado agregado, discutida
adiante, precisaria capturar.

\begin{table}[ht]
\centering
\caption{Taxa de detecção por tipo de ataque — IDS binário (granularidade por segundo)}
\label{tab:detection_rate}
\begin{tabular}{@{}lrrr@{}}
\toprule
\textbf{Ataque} & \textbf{Segundos de ataque} & \textbf{Segundos alertados} & \textbf{Taxa} \\ \midrule
recon             &  172 &  36 & 20,9\% \\
dos (incl.\ ddos) &  483 & 287 & 59,4\% \\
bruteforce        &  107 & 107 & \textbf{100,0\%} \\
web               &  101 & 101 & \textbf{100,0\%} \\
mitm              &  151 & 151 & \textbf{100,0\%} \\
spoofing          &  271 & 269 & \textbf{99,3\%} \\
malware           &   75 &  16 & 21,3\% \\ \bottomrule
\multicolumn{4}{l}{\small Janelas estendidas por \texttt{idle\_slack} = 90~s
                   (absorve o dreno do \textit{flood}). \textit{web} e \textit{mitm}} \\
\multicolumn{4}{l}{\small agora têm tráfego representativo (servidor HTTP real e
                   ARP \textit{spoofing} com tráfego forçado).}
\end{tabular}
\end{table}

\subsection{IDS Multiclasse (Fases 1+2)}
\label{subsec:res_multi}

Todas as oito classes geraram tráfego representativo no experimento: um servidor
HTTP real como alvo do ataque \texttt{web}, ARP \textit{spoofing} bidirecional
com tráfego forçado pela rota envenenada para \texttt{mitm} e \textit{flood} com
origem forjada para \texttt{spoofing}. Como as janelas de cada ataque foram
isoladas por 120~s de intervalo ocioso, os resultados por classe são
atribuíveis ao comportamento de cada ataque, e não a sobreposição temporal entre
janelas.

O desempenho por classe (Tabela~\ref{tab:multi_metrics}) divide os ataques em
dois regimes nítidos. A Fase~2 classifica corretamente \texttt{recon}
(F1 = 0,976) e parcialmente \texttt{malware} (F1 = 0,268), mas as classes
\texttt{bruteforce}, \texttt{web}, \texttt{mitm} e \texttt{spoofing} colapsam na
classe \texttt{dos}, que absorve quase todos os fluxos de ataque (revocação de
0,998 e precisão de apenas 0,46). Na prática, o rótulo de tipo é útil para
distinguir varreduras de reconhecimento, mas não separa esses quatro ataques de
uma inundação genérica, o que mantém a triagem dependente do alerta binário da
Fase~1 nesses casos.


\begin{table}[ht]
\centering
\caption{IDS multiclasse na validação no \textit{testbed}: 8 classes (ddos unificado a dos).}
\label{tab:multi_metrics}
\begin{tabular}{@{}lrrrrr@{}}
\toprule
\textbf{Classe} & \textbf{TP} & \textbf{FP} & \textbf{FN} & \textbf{Prec.} & \textbf{F1} \\ \midrule
recon             &    987 &     14 &     35 & 0,99 & \textbf{0,976} \\
dos (incl.\ ddos) & 13.302 & 15.468 &     29 & 0,46 & 0,632 \\
malware           &    373 &    137 &  1.898 & 0,73 & 0,268 \\
bruteforce        &      4 &     17 &  5.023 & 0,19 & 0,002 \\
web               &      0 &      9 &  2.737 & 0,00 & 0,000 \\
mitm              &      0 &      1 &  1.133 & 0,00 & 0,000 \\
spoofing          &      0 &      2 &  4.793 & 0,00 & 0,000 \\ \midrule
\textbf{Acurácia geral (tipo)} & \multicolumn{5}{r}{\textbf{48,4\%}} \\
\textbf{Macro F1} & \multicolumn{5}{r}{0,320} \\
\bottomrule
\end{tabular}
\end{table}


Esse colapso tem duas causas que se somam. A primeira é a limitação do próprio
espaço de \textit{features}: no vetor de 55 atributos por fluxo, ataques de
fluxos curtos e de alta taxa, como força bruta contra serviços que falham
rápido, varredura web e o encaminhamento ARP do MITM, tornam-se estatisticamente
indistinguíveis de uma inundação. Para \texttt{spoofing}, classificar como
\texttt{dos} é inclusive a resposta correta, já que se trata de um \textit{flood}
com origem forjada. Separar as demais exigiria estado agregado entre fluxos, como
contadores de falha de autenticação ou a diversidade de IPs de origem, ausentes
por construção das \textit{features} por fluxo. A segunda causa é o desvio entre
treino e produção: na avaliação \textit{offline}
(Tabela~\ref{tab:phase2_offline}) o modelo atinge F1 de 0,80, 0,53 e 0,31 para
\texttt{web}, \texttt{mitm} e \texttt{bruteforce}, mas o tráfego gerado ao vivo
por essas ferramentas difere do dos datasets CIC e recai em \texttt{dos}.
Avaliá-las de forma realista exigiria serviços e vítimas reais no ambiente de
implantação, limitação registrada nos trabalhos futuros.

\subsection{Desvio de Distribuição e Artefatos de Implantação}
\label{subsec:res_artefatos}

Dois fenômenos de implantação, ausentes de avaliações puramente \textit{offline},
ajudam a entender os resultados anteriores e são necessários para reproduzir o
sistema. O primeiro é o desvio de distribuição no tráfego benigno. Os datasets
CIC utilizam tráfego benigno \textit{scripted} com padrões temporais regulares,
enquanto o tráfego interativo real (SSH e HTTP de curta duração) apresenta
características distintas que o modelo classificava como ataque. As duas
\textit{features} de maior ganho respondem pela maior parte desse erro:
\texttt{flow\_iat\_min} está na faixa de milissegundos no treino benigno, mas em
\textit{bursts} TCP reais cai para 1--10~µs, de modo que qualquer \textit{burst}
no ambiente real parece anômalo ao modelo; do mesmo modo, fluxos curtos sobre
rede local rápida produzem \texttt{flow\_byts\_s} muito acima do observado no
treino. Embora tráfego benigno real já tenha sido incorporado ao treino
(Seção~\ref{sec:metodologia}), o desvio residual indica que a solução mais
completa é capturá-lo no próprio ambiente de implantação, a rede do VIM~4.

O segundo fenômeno é um artefato do \textit{netflower}: cada encerramento de
conexão TCP gera um micro-fluxo adicional com $\texttt{flow\_duration} = 0$,
$\texttt{tot\_fwd\_pkts} = 1$, $\texttt{totlen\_fwd\_pkts} = 52$~bytes e apenas
as \textit{flags} FIN e ACK. Como esse padrão não tem correspondente no treino
benigno, o modelo lhe atribui sistematicamente cerca de 84\% de confiança de
ataque, uma fonte de falsos positivos que independe do tráfego real. A mitigação
é uma guarda de pacote mínimo (\texttt{tot\_fwd\_pkts} $\geq 2$) no script do
IDS, que descarta esses micro-fluxos antes da inferência.

\subsection{Comparação entre IDS Binário e IDS Multiclasse}
\label{subsec:res_comp}
A Tabela~\ref{tab:comparison} consolida a comparação entre os dois pipelines
testados no VIM~4. As duas sessões usam o mesmo modelo de Fase~1 e o mesmo roteiro de ataques, com baseline benigno em ambas, por isso a detecção binária e o FPR são comuns aos dois pipelines (a Seção~\ref{subsec:res_bin} os mede de forma dedicada). Adicionalmente, como a VIM~4 não expõe interfaces de telemetria de energia em software (RAPL, \textit{Running Average Power Limit}, ou o sensor INA226), adota-se uma estimativa calibrada e não medição direta. O método e suas limitações são descritos ao final desta seção.

\begin{table}[ht]
\centering
\caption{Comparação dos dois pipelines avaliados no VIM~4: binário (Fase~1) vs.\ hierárquico completo (Fases~1+2).}
\label{tab:comparison}
\begin{tabular}{@{}lrr@{}}
\toprule
\textbf{Métrica} & \shortstack[r]{\textbf{IDS Binário}\\ \textbf{(só Fase~1)}}
                 & \shortstack[r]{\textbf{IDS Multiclasse}\\ \textbf{(Fases~1+2)}} \\ \midrule
\multicolumn{3}{l}{\textit{Detecção binária (Fase 1)}} \\
\quad Precisão (por seg.) & \textbf{89,3\%} & mesmo modelo P1 \\
\quad Revocação (por seg.) & 71,1\%         & mesmo modelo P1 \\
\quad F1-score            & 79,2\%          & mesmo modelo P1 \\
\quad FPR (baseline limpo) & \textbf{0,0\%} & \textbf{0,0\%} \\ \midrule
\multicolumn{3}{l}{\textit{Classificação de tipo (Fase 2)}} \\
\quad Acurácia (tipo)   & —               & \textbf{48,4\%} \\
\quad recon F1          & —               & \textbf{0,976} \\
\quad dos F1            & —               & 0,632 \\
\quad malware F1        & —               & 0,268 \\
\midrule
\multicolumn{3}{l}{\textit{Recursos (VIM 4, durante a sessão)}} \\
\quad CPU média         & 12,9\%          & 13,9\% \\
\quad CPU máxima        & 29,0\%          & 31,7\% \\
\quad RAM média         & 1.187 MB        & 1.238 MB (\textbf{+50 MB}) \\
\quad RAM máxima        & 1.263 MB        & 1.333 MB (+70 MB) \\ \midrule
\multicolumn{3}{l}{\textit{Throughput}} \\
\quad Alertas em ataque & \textbf{86.887}  & 30.314 \\
\quad Taxa              & \textbf{62,9 alertas/s} & 22,1 alertas/s \\ \midrule
\multicolumn{3}{l}{\textit{Energia estimada$^{b}$}} \\
\quad Potência média (ataque) & 3,46~W & 3,50~W \\
\quad Potência média (idle)   & 2,90~W & 3,00~W \\
\bottomrule
\end{tabular}
\end{table}





O resultado central da comparação é que o custo de acrescentar a Fase~2 ao nó de
borda é quase inteiramente de memória. Durante a sessão, o uso de CPU é
praticamente idêntico nos dois pipelines ($\approx$13\% em média), e o único
acréscimo mensurável do segundo modelo são 50~MB de RAM, cerca de 4\% sobre a
linha de base do binário. Isso ocorre porque a Fase~2 executa a inferência em
série após a Fase~1, mas sua latência é amortizada pelo rastreamento de fluxos,
etapa que já domina o tempo de processamento. O valor de RAM deriva das amostras
periódicas do log (\texttt{SYS\_SNAPSHOT}), e não de um único bloco de resumo, o
que lhe dá base sólida. Na prática, rotular o tipo de ataque para triagem é
viável no próprio dispositivo de borda sem comprometer seu orçamento de CPU ou
energia.

Esse ganho de informação vem, porém, com menor volume de alertas. O IDS binário
emite alerta para todo fluxo com $P_1 \geq 0{,}9$, a 62,9 alertas/s, enquanto o
multiclasse aplica um filtro adicional pelo limiar da Fase~2 e reduz o volume
bruto a 22,1 alertas/s, uma taxa 2,8$\times$ menor. A escolha entre os dois
pipelines é, portanto, orientada pelo uso: onde se precisa de alta frequência de
alertas brutos, por exemplo para alimentar um sistema de gestão de eventos e
informações de segurança (\textit{Security Information and Event Management},
SIEM), o binário é preferível; onde se precisa de rótulo de tipo para triagem, o
multiclasse agrega valor sem custo de CPU.

A comparação também confirma, agora lado a lado, os dois comportamentos de
detecção já discutidos: sobre o baseline benigno limpo, ambos os pipelines
mantêm FPR de 0\%, e o FNR de 28,9\% por segundo do IDS binário concentra-se em
\textit{recon}, \texttt{dos} e \textit{malware}, enquanto as classes de inundação
e curta duração são detectadas em pelo menos 99\% dos segundos.

Um achado operacional condiciona todos esses números: sob \textit{flood} de alta
taxa, o nó ARM não processa os fluxos em tempo real e o extrator os emite com
atraso, deixando um rastro de cerca de 80~s após o fim do ataque, a
aproximadamente 98 fluxos/s. Esse rastro inflava o FPR aparente e contaminava as
janelas seguintes. A mitigação, com 120~s de intervalo ocioso entre ataques e
\texttt{idle\_slack} de 90~s na análise, isola as janelas (apenas 14\% dos
alertas caem fora de janela) e foi o que permitiu obter resultados por classe não
contaminados. O fenômeno em si é um limite de capacidade do nó de borda sob
ataque volumétrico, a ser tratado, por limitação de taxa dos \textit{floods} ou
captura mais eficiente, em trabalhos futuros.

O consumo de energia da Tabela~\ref{tab:comparison} é uma estimativa, não uma
medição direta, porque o VIM~4 não expõe interface RAPL, INA226 ou
\texttt{hwmon}. A potência é modelada linearmente pela utilização de CPU,

\begin{equation}
P(t) = P_{idle} + (P_{max} - P_{idle}) \times \text{CPU\%}(t),
\label{eq:power}
\end{equation}

com \textit{endpoints} calibrados: \texttt{calibrate\_power.py} mede a utilização
de CPU em repouso (0,1\%) e sob carga total (\texttt{stress}, 99,9\%),
registrando como proveniência a temperatura (de 49 a 64~°C) e a frequência por
\textit{cluster} (A73 de 1800 a 2208~MHz; A53 de 500 a 2016~MHz), e ancora
$P_{idle} = 2{,}5$~W e $P_{max} = 9{,}0$~W no envelope de potência
\textit{board-level} do Khadas VIM~4 / A311D2 (USB-C PD). Os scripts do IDS
registram \texttt{power\_w}, temperatura e o escalonamento dinâmico de tensão e
frequência (\textit{Dynamic Voltage and Frequency Scaling}, DVFS) a cada amostra
($\sim$3~s), e a potência é integrada sobre essa série, separando ataque e
repouso pelas janelas de \textit{ground truth}. A incerteza é reportada como uma
banda de sensibilidade ($P_{idle} \in [2{,}0;3{,}0]$~W, $P_{max} \in
[7{,}0;11{,}0]$~W): de 4.312 a 6.576~J no binário e de 4.355 a 6.646~J no
multiclasse. O que importa é a comparação relativa: o consumo durante ataques é
praticamente o mesmo nos dois pipelines ($\approx$3,5~W), coerente com o fato de
o custo da Fase~2 ser de RAM, e não de CPU. Uma medição absoluta exigiria um
medidor de hardware externo, apontado como trabalho futuro.

%─────────────────────────────────────────────────────────────────────────────
\section{Conclusão}
\label{sec:conclusao}

Este trabalho apresentou um IDS hierárquico de duas fases para ambientes IoT em
redes 6G com MEC, cobrindo desde a construção do pipeline de dados até a
implantação e a validação no \textit{testbed} sobre a placa VIM~4. A contribuição
central não é apenas o modelo, mas a demonstração de que o pipeline completo é
executável dentro do orçamento de recursos de um nó de borda ARM, com a
delimitação empírica de onde as \textit{features} por fluxo deixam de ser
suficientes. Em validação com o limiar de implantação fixado em 0,9, o IDS
binário alcançou FPR de 0\% sobre o baseline benigno limpo, precisão de 89,3\% e
detecção de pelo menos 99\% dos ataques de força bruta, web, MITM e
\textit{spoofing}. Sobre os fluxos assim sinalizados, a Fase~2 classificou
corretamente o reconhecimento de rede (F1 = 0,976), a um custo adicional de
apenas 50~MB de RAM, sem \textit{overhead} de CPU nem de energia, o que confirma
a premissa que motiva a arquitetura: rotular o tipo de ataque para triagem é
viável no próprio nó de borda, pois a inferência da segunda fase é amortizada
pelo rastreamento de fluxos, etapa que já domina o tempo de processamento.

A validação também tornou explícito o teto da abordagem atual. A Fase~2 converge
para \texttt{dos} as classes cujo comportamento ao vivo é, de fato, uma
inundação de fluxos curtos e de alta taxa (\texttt{bruteforce}, \texttt{web},
\texttt{mitm} e \texttt{spoofing}), colapso que para o \texttt{spoofing} é
inclusive a resposta correta. Essa fronteira é um resultado, não uma falha do
pipeline: separar tais ataques exige \textit{features} de estado agregado entre
fluxos, como contadores de falha de autenticação e diversidade de IPs de origem,
que as 55 \textit{features} por fluxo não capturam. Somam-se a ela dois
fenômenos operacionais que só a validação com tráfego real revelou: o desvio
entre o tráfego de treino e o de produção, que induz falsos positivos
contornados por baseline benigno limpo e guarda de pacote mínimo, e a saturação
do extrator de fluxos sob \textit{flood}, um limite de capacidade do nó ARM que
precisou ser isolado por intervalos ociosos entre ataques. Toda a validação foi
consolidada em scripts reprodutíveis, executáveis com um único comando, e o
projeto completo, incluindo os modelos treinados, está disponível em repositório
aberto\footnote{\url{https://github.com/luiz-linkezio/A-Hierarchical-Solution-based-on-MEC-and-IoT-for-Cyberattack-Detection-in-6G-Networks}};
a ferramenta \textit{netflower} é mantida em repositório próprio
(Seção~\ref{sec:metodologia}).

Esses limites orientam os trabalhos futuros. O mais imediato é enriquecer o
espaço de \textit{features} com estado agregado entre fluxos, capaz de separar
força bruta, \textit{spoofing} e MITM da classe \texttt{dos}, acompanhado de um
ambiente de avaliação com serviços e vítimas reais que reduza o desvio entre
treino e produção. Em paralelo, cabe mitigar a saturação do extrator sob
\textit{flood}, por limitação de taxa ou captura mais eficiente, reduzir o
\texttt{idle\_timeout} para capturar ataques de curta duração como o \textit{scan}
UDP do \textit{recon}, e substituir a estimativa de energia por medição direta
com um medidor de hardware externo.

%─────────────────────────────────────────────────────────────────────────────
\bibliographystyle{sbc}
\bibliography{artigo}
\end{document}